% !TEX program = xelatex
% =====================================================================
%  Дипломын ажлын үндсэн файл (XeLaTeX-ээр хөрвүүлнэ)
%  Загвар: Asura-lab/diplom (М.Мөнхдорж, 2026) — удирдагч багшийн зөвлөснөөр
%  ашиглав. Thesis.cls, нүүр хуудсууд загвараас өөрчлөлтгүй.
%  Хөрвүүлэх: VS Code + LaTeX Workshop (XeLaTeX) эсвэл `latexmk` (latexmkrc бий).
% =====================================================================
\documentclass[11pt]{Thesis}
\usepackage{fontspec}
% Windows дээр Times New Roman байна. Linux/Overleaf дээр хэмжээгээр нь
% ижил Liberation фонт руу шилжинэ (зөвхөн шалгах үед).
\IfFontExistsTF{Times New Roman}{\setmainfont{Times New Roman}}{\setmainfont{Liberation Serif}}
\IfFontExistsTF{Arial}{\setsansfont{Arial}}{\setsansfont{Liberation Sans}}
\IfFontExistsTF{Consolas}{\setmonofont[Scale=0.9]{Consolas}}{\setmonofont[Scale=0.9]{DejaVu Sans Mono}}
\usepackage{polyglossia}
\setdefaultlanguage{mongolian}
\setotherlanguage{english}
\usepackage{wrapfig}
\usepackage{multirow}
\usepackage{ragged2e}
\usepackage{amssymb}
\usepackage{indentfirst}
\usepackage{multicol}
\usepackage{listings}
\usepackage{adjustbox}
\usepackage{tikz}
\usepackage{xcolor}
\usepackage{colortbl}
\usepackage{enumitem}
\usetikzlibrary{positioning,arrows.meta,shapes.geometric,fit,calc,backgrounds}

% ---------------------------------------------------------------------
%  Тэмдэглэл үлдээх бол мөрийн эхэнд % тавина. Тэр мөр PDF-д харагдахгүй.
% ---------------------------------------------------------------------

\renewcommand{\lstlistingname}{Код}
\lstset{
    basicstyle=\ttfamily\small,
    breaklines=true,
    frame=single,
    showstringspaces=false
}

\begin{document}
\frenchspacing
\captionsetup{labelfont=bf, textfont=it}
\pagenumbering{roman}
\pagestyle{fancy}

% Улаан хаалттай мэдээллийг өөрийнхөөрөө солино. Сэдвийн нэр ажлын нэр
% тул удирдагч багштай баталгаажуулна.
\newcommand{\universityname}{ШИНЭ МОНГОЛ ТЕХНОЛОГИЙН КОЛЛЕЖ}
\newcommand{\departmentname}{КОМПЬЮТЕРЫН УХААНЫ ТЭНХИМ}
\newcommand{\titlename}{ШИНЭ МОНГОЛ ТЕХНОЛОГИЙН КОЛЛЕЖИЙН AI ЧАТБОТ: МОНГОЛ ЯРИА ТАНИХ БА ЭХ СУРВАЛЖТАЙ ХАРИУЛТ}
\newcommand{\authorsname}{\textcolor{red}{[Овог НЭР]}}
\newcommand{\authorShort}{\textcolor{red}{[О.Нэр]}}
\newcommand{\studentID}{\textcolor{red}{[оюутны код]}}
\newcommand{\supervisorname}{\textcolor{red}{[Удирдагч багш]}}
\newcommand{\cityname}{Улаанбаатар хот}
\newcommand{\gradYear}{2027}


%---------  Нүүр хуудас / Дотор нүүр
\begin{titlepage}

\raggedright

\noindent\begin{minipage}{0.15\textwidth}
\vspace{-24mm}
\includegraphics[width=\linewidth]{Pictures/NMCT.png}
\end{minipage}
\hfill
\begin{minipage}{0.8\textwidth}
\centering
\vspace{-24mm}
\textbf{\textsc{\Large \universityname}}\\[1.2mm]
\textbf{\textsc{\Large \departmentname}}\\
\end{minipage}
\noindent

\vspace{18mm}
\raggedleft
\text{ Оюутны код: }\studentID\\
\text{ Оюутны овог нэр: }\authorsname\\

\vspace{48mm}
\centering
{\linespread{1.08}\selectfont\textbf{\Large \parbox{0.95\textwidth}{\centering \titlename}}\\[2.5mm]}
\normalsize /ДЭД БАКАЛАВРЫН ТӨГСӨЛТИЙН СУДАЛГААНЫ АЖИЛ/\\

\vspace{20mm}
\begin{center}
\begin{minipage}{0.25\textwidth}
Удирдагч багш\\[1.5mm]
Гүйцэтгэсэн оюутан
\end{minipage}
\hspace{0.5cm}
\begin{minipage}{0.3\textwidth}
\raggedleft
/\supervisorname, Магистр/\\[1.5mm]
/\studentID, \authorShort/
\end{minipage}
\end{center}

\vspace{75mm}

\begin{center}
 \cityname  \\
 \gradYear \hspace{1pt} он
\end{center}

\end{titlepage}

\begin{titlepage}
\begin{center}
\vspace{-45mm}
\textbf{\textsc{\Large \universityname}}\\[1.2mm]
\textbf{\textsc{\Large \departmentname}}\\

\vspace{75mm}

\normalsize Дэд бакалаврын төгсөлтийн судалгааны ажил (061301)\\[2.5mm]
{\linespread{1.08}\selectfont\textbf{\Large \parbox{0.95\textwidth}{\centering \titlename}}\\}


\vspace{20mm}
\begin{center}
\begin{minipage}{0.4\textwidth}
Удирдагч багш\\[1.5mm]
Гүйцэтгэсэн оюутан
\end{minipage}
\hspace{0.5cm}
\begin{minipage}{0.3\textwidth}
\raggedleft
/\supervisorname, Магистр/\\[1.5mm]
/\studentID, \authorShort/
\end{minipage}
\end{center}

\end{center}

\vspace{60mm}

\vfill
\begin{center}
 \cityname \hspace{1pt} \\
 \gradYear \hspace{1pt} он\\
\end{center}

\end{titlepage}


%---------  Хураангуй (хамгийн сүүлд бичнэ)
% \newpage
% \addtotoc{Хураангуй}

\vspace{10mm}

\begin{center}
\textbf{\large Хураангуй}
\end{center}

\sloppy

\bichne{Хураангуйг бүх үр дүн гарсны дараа шинэчилнэ. Англи хураангуйг (150--250 үг) урьдчилсан хамгаалалтын дараа нэмнэ.}

Энэ ажлаар Шинэ Монгол Технологийн Коллежийн албан мэдээллийн талаар Монгол хэлээр бичгээр болон дуугаар асууж, эх сурвалжтай хариулт авах AI чатбот хөгжүүлж байна. Систем ярьсан асуултыг текст болгоно. Дараа нь сургуулийн баримтаас тохирох хэсгийг хайж, олсон хэсгээ дүрмээр сонгон эх сурвалжтай нь харуулна. Хариултыг дуугаар уншиж өгнө.

Монгол хэлний дуу таних хоёр аргыг нэг хэлэгчийн ижил 12 бичлэг дээр харьцуулсан. Нэг нь Монгол хэлээр нэмэлт сургасан Whisper small загвар, нөгөө нь Google-ийн Gemini 3.5 Transcribe үйлчилгээ. Үгийн алдааны хувь (WER) Whisper-т 38.07\%, Gemini-д 40.34\% гарсан. Тэмдэгтийн алдааны хувь (CER) 17.93\% ба 24.85\% байв. Gemini нэг бичлэгт хэлсэн оныг буруу таньсан. Ялгаа бага, бичлэг цөөн тул аль арга нь ерөнхийдөө илүү гэж дүгнэх боломжгүй. Иймээс сонголтыг хэмжсэн алдаа, үйлчилгээний зориулалт, зардал, интернэт ба өгөгдлийн аюулгүй байдал, Монгол хэлэнд сайжруулах боломжийг хамтад нь харгалзан хийнэ. Хайлт болон хариулт сонголтын үнэлгээг урьдчилан түгжсэн асуултын багцаар цаашид хийнэ.

\vspace{3mm}
\textit{Түлхүүр үгс:} Монгол хэлний яриа таних, Whisper, Gemini, RAG, эх сурвалжтай хариулт, чатбот


%---------  Гарчиг
\newpage
\tableofcontents

%---------  Товчилсон үгсийн жагсаалт
\addtotoc{Товчилсон үгсийн жагсаалт}
\chapter*{Товчилсон үгсийн жагсаалт}

% Шинэ товчлол хэрэглэх бүрдээ энд \acro{...}{...} нэмнэ.
\begin{acronym}[ASR]
\acro{API}{Application Programming Interface -- Програмын харилцах гүүр}
\acro{ASR}{Automatic Speech Recognition -- Яриа автоматаар таних}
\acro{CER}{Character Error Rate -- Тэмдэгтийн алдааны хувь}
\acro{GPU}{Graphics Processing Unit -- График процессор}
\acro{RAG}{Retrieval-Augmented Generation -- Хайлтаар баяжуулсан хариулт}
\acro{RRF}{Reciprocal Rank Fusion -- Байрын урвуугаар нэгтгэх}
\acro{TTS}{Text-to-Speech -- Текстийг дуу болгох}
\acro{WER}{Word Error Rate -- Үгийн алдааны хувь}
\end{acronym}


%---------  Хүснэгт, зургийн жагсаалт
\newpage
\listoftab
\newpage
\listoffig

%---------  Үндсэн хэсэг
\newpage
\pagenumbering{arabic}
\setcounter{page}{1}

\chapter*{Ажлын төлөвлөгөө}
\addcontentsline{toc}{chapter}{Ажлын төлөвлөгөө}
\label{ch:plan}

\noindent\textbf{Оюутны нэр:} \authorShort \hfill \textbf{Удирдагч багшийн нэр:} \supervisorname

\noindent\textbf{Судалгааны ажлын сэдэв:} \titlename % сэдвийн нэрийг удирдагчтай баталгаажуулна

\vspace{0.2cm}

Ажлыг долоо хоног бүрээр төлөвлөсөн. Долоо хоног бүр нэг гол ажил хийж, тухайн ажилтай холбоотой нэг судалгааны өгүүлэл уншина. Улаанаар тэмдэглэсэн нь сургуулийн гарын авлагад заасан үзлэг, хамгаалалт юм.

{\small
\setlength{\tabcolsep}{4pt}
\renewcommand{\arraystretch}{1.12}
\begin{longtable}{|p{1.0cm}|p{0.9cm}|p{1.85cm}|p{6.0cm}|p{3.0cm}|p{1.3cm}|}
\caption{Ажлын долоо хоногийн төлөвлөгөө (2026.10--2027.05)}\label{tab:plan}\\
\hline
\textbf{Сар} & \textbf{7 хоног} & \textbf{Огноо} & \textbf{Гол ажил} & \textbf{Судалгаа (унших)} & \textbf{Гүйцэт\-гэл} \\
\hline
\endfirsthead
\hline
\textbf{Сар} & \textbf{7 хоног} & \textbf{Огноо} & \textbf{Гол ажил} & \textbf{Судалгаа (унших)} & \textbf{Гүйцэт\-гэл} \\
\hline
\endhead
10 сар & W01 & 10.26--11.01 & Ажлаа цэгцэлж, системийн гол сонголтуудыг хийх & -- & $\checkmark$ \\
\hline
11 сар & W02 & 11.02--11.08 & Семинар хийж, багштай тодруулах; Whisper--Gemini харьцуулалтыг дипломд бичих & RAG & Явж байна \\
 & W03 & 11.09--11.15 & Дуу таних аргаа шинэ бичлэгээр харьцуулж сонгох & Хариугүй асуулт (SQuAD 2.0) & \\
 & W04 & 11.16--11.22 & Шалгах асуултын багц бэлдэж түгжих & BM25 & \\
 & W05 & 11.23--11.29 & Үзлэг №1-ийн материал бэлдэх & RRF, тойм бичих & \\
\hline
12 сар & W06 & 11.30--12.06 & \textcolor{red}{Үзлэг №1 (12/01)} & Давтлага & \\
 & W07 & 12.07--12.13 & Хайлтын аргуудыг харьцуулах (мэдлэгийн модтой, модгүй) & Multilingual E5 & \\
 & W08 & 12.14--12.20 & Хариулт сонгох хоёр аргыг (дүрэм, LLM) харьцуулах & Эшлэлтэй хариулт & \\
 & W09 & 12.21--12.27 & Мэдлэгийн модыг баяжуулах: багш, өрөө, анги, хуваарь & RAPTOR, GraphRAG & \\
 & W10 & 12.28--01.03 & Завсрын дүгнэлт хийх & Давтлага & \\
\hline
1 сар & W11 & 01.04--01.10 & Нэр, тоо, оныг дуугаар зөв уншуулах & Тоо хэвийн болгох & \\
 & W12 & 01.11--01.17 & Хариултыг сонсож ойлгосныг шалгах & TTS-ийн ойлгомж & \\
 & W13 & 01.18--01.24 & Хариулт ба дуу танилтын гол алдааг засах & Дуу танилтын алдаа ба нэр & \\
 & W14 & 01.25--01.31 & Мэдээлэл шинэчлэх, тасалдлыг шалгах & DPR & \\
\hline
2 сар & W15 & 02.01--02.07 & Бичсэн ба ярьсан асуултыг харьцуулах; \textcolor{red}{Үзлэг №2} & Утгын хэмжүүр & \\
 & W16 & 02.08--02.14 & Өөр компьютер дээр ажиллуулах & Хэрэглэгчийн хүлээлт & \\
 & W17 & 02.15--02.21 & Эцсийн хувилбар, LLM-ийн эцсийн шийдвэр & LLM ба RAG & \\
 & W18 & 02.22--02.28 & Гол хэмжилтүүдийг хийх & RAGAS & \\
\hline
3 сар & W19 & 03.01--03.07 & Хүний үнэлгээ хийх & TTS-ийн хүний үнэлгээ & \\
 & W20 & 03.08--03.14 & Үр дүн, алдааны шинжилгээг бичих & Bootstrap статистик & \\
 & W21 & 03.15--03.21 & \textcolor{red}{Бүтэн ноорог (40 хүртэл хуудас)} & Давтлага & \\
 & W22 & 03.22--03.28 & Удирдагчийн саналаар засах & -- & \\
 & W23 & 03.29--04.04 & Үзүүлэн, хамгаалалтын яриа бэлдэх & -- & \\
\hline
4 сар & W24 & 04.05--04.11 & Урьдчилсан хамгаалалтын материал & -- & \\
 & W25 & 04.12--04.18 & \textcolor{red}{Урьдчилсан хамгаалалт} & -- & \\
 & W26 & 04.19--04.25 & Комиссын саналаар засах & -- & \\
 & W27 & 04.26--05.02 & Засварыг дуусгах, давтан ажиллуулах заавар & -- & \\
\hline
5 сар & W28 & 05.03--05.09 & Эцсийн бичвэр, Монгол ба англи хураангуй & -- & \\
 & W29 & 05.10--05.16 & Товчлол (A4, 2--4 хуудас), хүлээлгэн өгөх материал & -- & \\
 & W30 & 05.17--05.23 & Хамгаалалтын давтлага & -- & \\
 & W31 & 05.24--05.30 & \textcolor{red}{Эцсийн хяналт, хүлээлгэн өгөх} & -- & \\
 & W32 & 05.31 & Нөөц өдөр & -- & \\
\hline
\end{longtable}
}

% Үзлэг №2, урьдчилсан ба эцсийн хамгаалалтын яг огноог сургууль зарлахаар шинэчилнэ.
                    % Ажлын төлөвлөгөө
\chapter{Удиртгал}
\label{ch:introduction}

%======================================================================
% 2026.11.11 (3): бүтэц, гарчгийг сургуулийн өмнөх дипломтой (Asura-lab/diplom)
% ижил болгосон: 1.1 Үндэслэл, ач холбогдол; 1.2 Зорилго, зорилт (дэд хэсэгтэй);
% 1.3 Судалгааны асуулт ба хэмжигдэх шалгуур. Үг хэллэгийг Монгол 7 дипломын
% удиртгалтай тулгасан (papers/28-mn-thesis-style-20261111.mn.md). Агуулга хуулаагүй.
% 2026.11.11 (2): үндэслэлийг баримтаар, зорилтыг судалгааны ажлын хэлбэрээр бичсэн;
% дуу таних жижиг хэмжилтийг зорилт, судалгааны асуултаас хассан (3.6, 4.1-д үлдсэн).
% Өмнөх бүтэн хувилбар: thesis/archive-20261111/.
%======================================================================

\section{Үндэслэл, ач холбогдол}
\label{sec:intro-background}

Шинэ Монгол Технологийн Коллежийн тухай мэдээлэл нэг газарт төвлөрөөгүй. Элсэлтийн журам, сургалтын төлбөр, хөтөлбөрийн танилцуулга сургуулийн вэб сайтад бий. Харин аль өрөө хэдэн давхарт байдаг, аль анги хэдэн цагт хаана хичээллэдэг, ямар асуудлаар аль албаны хэнд хандах вэ гэдэг нь сайтад байхгүй. Эдгээр нь сургалтын албаны өрөөний жагсаалт, хичээлийн хуваарь, ажилтны жагсаалтад хадгалагддаг. Иймээс сургуульд ирсэн оюутан, элсэгч, эцэг эх, зочин ийм мэдээллийг ажилтнаас биечлэн асууж авдаг.
% ӨӨРИЙН ҮГЭЭР: ресепшнээс цуглуулсан бодит асуултын 2–3 жишээ, аль төрлийн асуулт хамгийн
% их давтагддагийг энд нэм (боломжтой бол хэдэн хүнээс, хэдийд асуусныг хүснэгтээр).
% Агент асуулт, тоо зохиохгүй.

Тулгамдаж буй асуудлыг тодруулахын тулд 2026 оны 11-р сард сургалтын албаны ажилтнуудтай уулзав. Ажилтнууд хоёр зүйлийг онцолсон. Нэгдүгээрт, тэд ижил мэдээллийг өөр өөр хүнд дахин дахин тайлбарладаг. Хоёрдугаарт, нэгдүгээр давхарт цаасан мэдээлэл байрлуулсан боловч хүмүүс түүнээс хариултаа олохгүй, ажилтнаас дэлгэрүүлж асуудаг. Мөн ийм систем зөвхөн ресепшнд биш, гар утаснаас нээгддэг байвал илүү олон хүнд хүрнэ гэж зөвлөсөн.
% ӨӨРИЙН ҮГЭЭР: уулзсан өдөр, хэнтэй (албан тушаал) уулзсаныг нэм.

Ийм бэрхшээл бусад сургуульд ч ажиглагдсан. Тухайлбал, Милано-Бикокка их сургуулийн оюутнууд анги танхим, хичээлийн хуваарь, шалгалтын өрөөний мэдээллийг олоход хэцүү гэж хариулсан \cite{antico2024unimib}. Burch их сургуулийн оюутнууд мэдээллээ сайтаас олж чадаагүй үедээ энгийн асуултаар захиргааны ажилтанд ханддаг бөгөөд энэ нь ажилтны ачааллыг нэмдэг \cite{matic2026burch}. Эдгээр сургууль асуудлыг оюутны асуултад сургуулийн баримтаас хайж хариулдаг туслах системээр шийдэхийг зорьсон.

Ийм туслахыг Монгол хэл дээр хийхэд хоёр бэрхшээл тулгардаг. Нэгдүгээрт, Монгол хэл яриа танилтын нөөц багатай хэлэнд тооцогддог. Олон хэлний Whisper загвар Монгол яриаг 100\%-иас их үгийн алдаатай таньсан \cite{radford2023whisper}, Халх Монгол хэлний нээлттэй сургалтын яриа хэдхэн цаг л бий \cite{mengke2025khalkha}. Хоёрдугаарт, сургуулийн нэрийн өмнөөс өгөх хариулт үнэн зөв байх ёстой. Хэлний загвар байхгүй баримтыг зохиож болдог тул хариултыг баримтын сангаас хайж, олсон эхдээ тулгуурлан өгөх шаардлагатай \cite{lewis2020rag}.

Иймд сургуулийн тархай бутархай мэдээллийг нэг бүтцэд оруулж, Монгол хэлээр асуусан асуултад түүнээс хариулдаг, хариулт байхгүй үед хэнд, хаана хандахыг зааж өгдөг туслах систем боловсруулах шаардлага тодорхой болов.

Энэ ажлын практик ач холбогдол нь сургуулийн ажилтнуудын ижил мэдээллийг давтан тайлбарлах ачааллыг бууруулж, сургуульд ирсэн хүн хэрэгтэй мэдээллээ нэг газраас авах боломжийг бүрдүүлэхэд оршино. Судалгааны ач холбогдол нь гараар зохион байгуулсан мэдлэгийн модыг хайлтад ашиглах аргыг Монгол хэл дээр, сургуулийн бодит мэдээллээр шалгахад оршино.

\section{Зорилго, зорилт}
\label{sec:intro-goals}

\subsection{Судалгааны зорилго}

Энэ судалгааны зорилго нь сургуулийн ресепшнд ирсэн хүний Монгол хэлээр ярьж эсвэл бичиж асуусан асуултад \textbf{сургуулийн албан мэдээлэлд тулгуурлан хариулдаг туслах системийг боловсруулахад} оршино.

\subsection{Судалгааны зорилтууд}

Энэ зорилгод хүрэхийн тулд дөрвөн зорилтыг дэвшүүлэв:
\begin{enumerate}
    \item Ресепшнд ирдэг хүмүүсийн мэдээллийн хэрэгцээ, системд тавигдах шаардлагыг сургуулийн албадтай уулзаж тодорхойлох.
    \item Сургуулийн мэдээллийг хоёр эхээс, тухайлбал албан вэб сайт болон сургуулийн албадаас цуглуулж, мэдлэгийн мод хэлбэрээр зохион байгуулах.
    \item Мэдлэгийн модонд тулгуурласан хайлтын аргыг боловсруулж, бодит асуултаар үнэлэх.
    \item Дуугаар болон бичгээр асуусан асуултад эх сурвалжтай хариулт өгөх системийг хэрэгжүүлэх.
\end{enumerate}
% Зорилт 1: карьерын төв, ресепшний ажилтантай уулзсаны дараа «сургуулийн албад»-ыг нэрээр нь бич.

\section{Судалгааны асуулт ба хэмжигдэх шалгуур}
\label{sec:rq}

Гуравдугаар зорилтыг хэмжихүйц шалгууртай холбохын тулд дараах судалгааны асуултыг тодорхойлов: асуултын сэдвийг мэдлэгийн модны хүрээгээр тодорхойлж, тухайн хүрээний эхэд давуу эрх өгөх нь бүх эхээс ялгалгүй хайхтай харьцуулахад зөв эхийг олох чадварыг сайжруулах уу?

Хэмжих шалгуур нь зөв эх хайлтын эхний нэг болон эхний гурван үр дүнд орсон асуултын тоо юм. Хоёр аргыг ижил асуултаар харьцуулна. Туршилтын аргыг \ref{sec:exp-tree}-т тодорхойлов.

\section{Судалгааны хамрах хүрээ}
\label{sec:intro-scope}

Систем зөвхөн сургуулийн албан мэдээллээр хариулна. Мэдээллийг хоёр эхээс авсан: сургуулийн вэб сайтын нийтэд нээлттэй хуудсууд болон сургуулийн албадаас өгсөн өрөө, анги бүлэг, ажилтан, хичээлийн хуваарийн жагсаалт. Асуулт, хариулт Монгол хэлээр байна. Дүн, ирц зэрэг хувь хүний мэдээлэл болон ажилтны хувийн утасны дугаарыг энэ ажилд хамруулаагүй.

\section{Судалгааны ажлын бүтэц}
\label{sec:intro-structure}

Энэ ажил таван бүлгээс бүрдэнэ. \ref{ch:literature}-р бүлэгт Монгол хэлний яриа танилт, хайлтаар баяжуулсан хариулт, хариулт байхгүй асуулт, мэдлэгийн мод дээрх хайлтын талаарх өмнөх судалгааг тоймлов. \ref{ch:method}-р бүлэгт системд тавигдах шаардлага, мэдээлэл цуглуулсан арга, системийн бүтэц, туршилтын аргыг тодорхойлов. \ref{ch:results}-р бүлэгт одоогоор гарсан үр дүнг, \ref{ch:conclusion}-р бүлэгт дүгнэлт, хязгаарлалт, цаашдын ажлыг тусгав.
   % 1. Удиртгал
\chapter{Судалгааны сэдвийн онол, өнөөгийн түвшин}
\label{ch:literature}

Энэ бүлэгт эхлээд Монгол хэлний яриа таних аргуудыг тоймлоно (§\ref{sec:lit-asr}). Дараа нь сургуулийн баримтаас хариулт хайх аргуудыг авч үзнэ (§\ref{sec:lit-rag}).

%======================================================================
\section{Монгол хэлний яриа таних}
\label{sec:lit-asr}

\subsection{Яриа таних ба алдааны хэмжүүр}

Яриа автоматаар таних (\ac{ASR}, \ac{STT}) нь хүний ярианы аудиог текст болгох ажил юм. Чатботод хүн асуултаа ярихад эхлээд энэ алхам ажиллана. Энд гарсан алдаа дараагийн бүх алхамд дамждаг. Жишээ нь «хураамж» гэдэг үгийг буруу таньвал систем өөр мэдээлэл хайж эхэлнэ.

Яриа танилтын чанарыг ихэвчлэн үгийн алдааны хувиар (\ac{WER}) хэмждэг:
\begin{equation}
\mathrm{WER} = \frac{S + D + I}{N} \times 100\%
\label{eq:wer}
\end{equation}
Энд $S$ нь өөр үгээр сольсон үгийн тоо, $D$ нь орхисон үгийн тоо, $I$ нь илүү нэмсэн үгийн тоо, $N$ нь зөв бичвэрийн (reference) нийт үгийн тоо. Тэмдэгтийн алдааны хувь (\ac{CER}) ижил томьёогоор үсгийн түвшинд тооцогдоно. Нэмсэн үгийн тоо их байвал WER 100\%-иас давж болно.

WER дангаараа утгын алдааг бүрэн харуулдаггүй. Жишээ нь зөв бичвэр нь «хоёр мянга хорин тав», таньсан текст нь «2025» байлаа гэж үзье (зохиомол жишээ). Утга нь ижил. Гэвч үгээр тоолоход зөв бичвэрийн 4 үгээс нэгийг сольж, гурвыг орхисон тул WER $= 4/4 = 100\%$ болно. Эсрэгээр «2026» гэснийг «2020» гэж таньвал ганц үгийн алдаа. Гэвч асуултын утга бүрэн өөрчлөгдөнө. Иймээс нэр, тоо, оны алдааг тусад нь шалгах шаардлагатай (§\ref{subsec:lit-eval}).

\subsection{Whisper загвар}

Radford нар (2023) Whisper загварыг интернэтээс цуглуулсан 680 мянган цагийн аудио ба бичвэрийн хосоор сургасан \cite{radford2023whisper}. Үүний 117 мянган цаг нь англиас бусад 96 хэлний яриа юм. Хүн гараар шошголоогүй, интернэт дэх бичвэрийг шууд ашигласан тул энэ аргыг сул хяналттай сургалт (weak supervision) гэдэг. Загвар нь encoder--decoder бүтэцтэй Transformer бөгөөд нэг загвар яриа таних, хэл таних, англи руу орчуулах зэрэг хэд хэдэн ажил хийдэг. Гол давуу тал нь тухайн өгөгдөлд нэмэлт сургалт хийлгүйгээр (zero-shot) шууд хэрэглэж болдог явдал.

Гэвч сургалтын өгөгдөлд Монгол яриа маш цөөн байсан. Иймээс том хэмжээний Whisper large-v2 загварын Монгол хэлний WER нь FLEURS өгөгдөл дээр 110.5\%, Common Voice 9 өгөгдөл дээр 128.4\% гарсан \cite{radford2023whisper}. \shalgah{110.5 ба 128.4-ийг өгүүллийн хавсралтын хүснэгтээс тулгах} Өөрөөр хэлбэл алдааны тоо зөв бичвэрийн үгийн тооноос ч их байна. Зохиогчид урт аудиод үг давтах, орхих, огт хэлээгүй бичвэр зохиох алдаа гардгийг хязгаарлалт болгон дурдсан \cite{radford2023whisper}.

Whisper нь tiny-ээс large хүртэл хэд хэдэн хэмжээтэй. Энэ ажилд ашигласан small загвар 244 сая параметртэй \cite{openai_whisper_repo}.

\subsection{Whisper-ийг Монгол хэлэнд нэмэлт сургах}

Үндсэн Whisper загварын Монгол хэлний алдаа өндөр тул Монгол өгөгдлөөр нэмэлт сургасан (fine-tuned) загварууд гарсан. Тэдний нэг нь Hugging Face дээр нээлттэй байршуулсан \texttt{Cafet/whisper-small-mongolian} загвар юм \cite{cafet_whisper_small_mn}. Энэ нь Whisper small-ийг Монгол хэлээр нэмэлт сургасан хувилбар. Загварын тайлбар хуудсанд сургалтын тохиргоо, шалгалтын WER бичигдсэн. Гэхдээ энэ нь рецензтэй судалгааны өгүүлэл биш. Мөн тэнд бичсэн WER манай сургуулийн асуултад ижил байна гэсэн баталгаа биш.

Mengke нар (2025) Халх Монгол хэлний өгөгдөл цөөн байх асуудлыг шийдэх аргыг санал болгосон \cite{mengke2025khalkha}. Common Voice 17-д Халх Монгол хэлний сургалтын өгөгдөл ердөө 3 цаг байдаг. Иймээс зохиогчид Өвөр Монголын Чахар аялгууны 147 цагийн өгөгдлийн бичвэрийг кирилл үсэг рүү хөрвүүлсэн (transliteration). Эхлээд энэ өгөгдлөөр, дараа нь Халхын өгөгдлөөр хоёр шаттай нэмэлт сургалт хийсэн. Whisper large-v3 загварын WER нь шууд хэрэглэхэд 92.60\%, зөвхөн Халхын өгөгдлөөр нэмэлт сургахад 43.02\%, хоёр шаттай аргаар 31.45\% болж буурсан. Хоёр өгөгдлийг шууд нийлүүлж сургахад 49.95\% гарсан. Сургалтад 48 GB санах ойтой A6000 GPU ашигласан \cite{mengke2025khalkha}.\footnote{Зохиогчид §3.2-т Whisper-ийг CTC алдагдлын функцээр сургасан гэж бичсэн. Энэ нь Whisper-ийн ердийн encoder--decoder сургалтаас өөр тул энэ хэсгийг нягтлах шаардлагатай. \shalgah{өгүүллийн §3.2-ыг дахин уншиж тулгах}}

Энэ судалгаанаас гурван сургамж гарна. Нэгдүгээрт, Whisper-ийг Монгол хэлэнд мэдэгдэхүйц сайжруулах боломжтой. Хоёрдугаарт, үүнд олон цагийн аудио--бичвэрийн хос, хүчтэй GPU, тусдаа шалгах өгөгдөл шаардлагатай. Гуравдугаарт, тэдний шалгалт уншсан яриа (Common Voice) дээр хийгдсэн тул сургуулийн чөлөөт асуултад ижил дүн гарна гэж хэлэх боломжгүй.

\subsection{Gemini 3.5 Transcribe}

Gemini 3.5 Transcribe нь Google-ийн Gemini \ac{API}-аар ажилладаг, аудио файлыг текст болгох зориулалттай үүлэн үйлчилгээ юм \cite{google_transcribe_model}. Албан баримтад Монгол хэл (\texttt{mn-MN}) дэмжих хэлний жагсаалтад орсон. Хэрэглэгч хэлийг урьдчилан зааж (language hint), үгчилсэн (verbatim) горимоор бичүүлж болно \cite{google_transcribe_api}.

Энэ үйлчилгээг ашиглахад \ac{API} түлхүүр, интернэт холболт шаардлагатай. Аудио Google-ийн сервер рүү илгээгддэг. Үнэгүй түвшинд (free tier) илгээсэн өгөгдлийг Google бүтээгдэхүүнээ сайжруулахад ашиглаж, хүн хянаж болно гэж нөхцөлд заасан \cite{google_terms}. Иймээс хувийн мэдээлэлтэй аудиог илгээхэд болгоомжлох хэрэгтэй. Төлбөртэй түвшний ойролцоо үнэ минутад 0.005 ам.доллар орчим байна (2026.10.29-ний байдлаар) \cite{google_pricing}.

Google-ийн өөр нэг загвар \texttt{gemini-3.5-live-translate-preview} нь яриаг өөр хэл рүү шууд орчуулах зориулалттай \cite{google_live_translate}. Түүний оролтын Монгол текстийг авч болох ч үндсэн зорилго нь орчуулга. Иймээс энэ ажилд түүнийг Transcribe-аас тусад нь авч үзсэн (§\ref{sec:res-demo}).

Gemini 3.5 Transcribe-ийн Монгол хэлний алдааг хэмжсэн рецензтэй судалгаа энэ ажлын хүрээнд олдоогүй. Баримтад хэл «дэмждэг» гэж бичсэн нь манай ярианд сайн таних баталгаа биш. Иймээс өөрсдийн бичлэг дээр хэмжих шаардлагатай болсон.

\subsection{Хоёр аргын харьцуулсан тойм}

\tref{tab:asr_compare}-т хоёр аргыг онолын талаас нь харьцуулав. Хамгийн гол ялгаа нь Whisper-ийг сайжруулахын тулд нэмэлт сургалт хийх шаардлагатай байхад Gemini-г API түлхүүрээр шууд ашиглаж болох явдал юм. Харин Gemini-д аудио гадагш гардаг, загварыг өөрөө сургах боломжгүй.

\begin{table}[H]
\centering
\caption{Whisper ба Gemini 3.5 Transcribe-ийн шинж чанарын харьцуулалт}
\label{tab:asr_compare}
\small
\begin{tabular}{|p{3.4cm}|p{5.1cm}|p{5.1cm}|}
\hline
\textbf{Шинж} & \textbf{Whisper (Cafet small)} & \textbf{Gemini 3.5 Transcribe} \\
\hline
Хаана ажиллах & Өөрийн компьютер (CPU эсвэл GPU) & Google-ийн үүлэн үйлчилгээ \\
\hline
Интернэт & Шаардлагагүй & Шаардлагатай \\
\hline
Аудио хаашаа очих & Компьютер дотроо үлдэнэ & Google-ийн сервер рүү очно \\
\hline
Монгол хэлэнд сайжруулах & Нэмэлт сургалт: олон цагийн өгөгдөл, GPU \cite{mengke2025khalkha} & Загварыг сургах боломжгүй; хэл заах, горим сонгох \\
\hline
Ашиглахад хэрэгтэй & Загварын файл, тооцоолох нөөц & API түлхүүр \\
\hline
Зардал & Тоног төхөөрөмж, цахилгаан & Үнэгүй квоттой; төлбөртэй нь минутад $\approx$0.005 ам.доллар \\
\hline
Нийтлэгдсэн Монгол хэлний хэмжилт & large-v2: FLEURS дээр WER 110.5\% \cite{radford2023whisper}; нэмэлт сургасан large-v3: 31.45\% \cite{mengke2025khalkha} & Рецензтэй хэмжилт олдоогүй \\
\hline
\end{tabular}
\end{table}

\subsection{Яриа танилтыг үнэлэхэд анхаарах зүйл}
\label{subsec:lit-eval}

Өмнөх судалгаанаас яриа танилтыг үнэлэхэд дөрвөн анхаарах зүйл гарсан.

\begin{enumerate}
    \item \textbf{Нэг арга бүх хэлэнд хамгийн сайн биш.} Jimerson нар 11 хэл дээр хэд хэдэн архитектурыг харьцуулахад нэг ч арга бүх хэлэнд тэргүүлээгүй \cite{jimerson2023unhelpful}\unshih. Иймээс сонголтоо өөрийн өгөгдөл дээр хийх хэрэгтэй.
    \item \textbf{Текстийг хэвийн болгох (normalization) нь дүнг өөрчилнө.} Цэг тэмдэг, тэмдэгтийг арилгах дүрэм зарим хэлэнд WER-ийг хиймлээр бууруулдаг \cite{manohar2024normalization}\unshih. Иймээс хоёр аргыг яг ижил дүрмээр тоолох ёстой.
    \item \textbf{Нэр, тооны алдаа дараагийн ажлыг эвддэг.} Яриа танилтын алдаа нэр томьёог таних ажилд хүчтэй нөлөөлдөг \cite{szymanski2023ner}\unshih. Манай чатботод нэр, тоо, он хайлтын гол түлхүүр болно.
    \item \textbf{Яриагүй хэсэгт бичвэр зохиох.} Whisper заримдаа огт хэлээгүй өгүүлбэр зохиодог \cite{koenecke2024careless}\unshih. Иймээс чимээгүй бичлэгээр шалгах хэрэгтэй.
\end{enumerate}

%======================================================================
\section{Мэдээлэл хайх ба RAG}
\label{sec:lit-rag}

\subsection{RAG-ийн үндсэн санаа}

Lewis нар (2020) \ac{RAG} аргыг санал болгосон \cite{lewis2020rag}\unshih. Энэ арга хоёр хэсгээс бүрдэнэ. Эхний хэсэг нь хайгч (retriever). Хайгч асуултад тохирох баримтын хэсгийг том сангаас олно. Зохиогчид Wikipedia-г 100 үгтэй 21 сая орчим хэсэг болгон хуваасан. Хоёр дахь хэсэг нь хариулт бичигч (generator). Тэр асуулт болон олсон хэсгийг уншиж хариулт бичнэ.

Энэ аргын гол санаа нь мэдлэгийг загварын параметр дотор биш, гадаад баримтын санд хадгалах явдал юм. Иймээс хоёр давуу тал гарна. Нэгдүгээрт, хариулт аль баримтаас гарсныг харуулах боломжтой. Хоёрдугаарт, загварыг дахин сургахгүйгээр баримтын санг шинэчилж мэдлэгээ өөрчилж болно. Сургуулийн мэдээлэл жил бүр өөрчлөгддөг тул энэ шинж манай ажилд чухал.

Жишээ авъя. Хэрэглэгч «Бүртгэлийн хураамж хэд вэ?» гэж асуув. Хайгч сургуулийн баримтаас бүртгэлийн хураамжийн тухай хэсгийг олно. Дараа нь хариулт тэр хэсэгт тулгуурлан гарна. Хэрэглэгч хариултын хажууд эх баримтыг нь харна.

Манай системд хайх хэсэг (R) нь гуравдугаар алхам буюу BM25 ба E5 хайлт юм (§\ref{sec:architecture}). Хариулт гаргах хэсэг (AG) нь дөрөвдүгээр алхам. Гэхдээ манай систем шинэ өгүүлбэр бичихгүй. Олсон хэсгүүдээс тохирох 1--3 хэсгийг дүрмээр сонгоод, эхийн үгээр нь харуулдаг. Давуу тал нь эх баримтад байхгүй зүйл зохиогдохгүй. Сул тал нь хэд хэдэн хэсгийн мэдээллийг нэгтгэж товч хариулж чаддаггүй. Хариулт сонгоход LLM ашиглах эсэхийг W08-д ижил асуултаар харьцуулна.

\bichne{Өгүүллийг уншсаны дараа энэ хэсгийг өөрийн үгээр засаж, RAG-Sequence ба RAG-Token хувилбарын ялгааг нэмнэ.}

\subsection{Үгээр хайх: BM25}
\bichne{W04 --- \cite{robertson2009bm25}\unshih: ховор үг яагаад их оноо авдаг, давтагдсан үгийн оноо яагаад хязгаартай вэ.}

\subsection{Утгаар хайх: вектор ба E5}
\bichne{W07 --- \cite{wang2024me5}\unshih: текстийг вектор болгох, косинусын ойролцоо, Монгол хэлний хязгаар.}

\subsection{Хоёр хайлтыг нэгтгэх: RRF}
\bichne{W05 --- \cite{cormack2009rrf}\unshih: $\sum 1/(k+\text{байр})$ томьёо; өгүүлэлд $k=60$, манайд $k=30$.}

\subsection{Мэдлэгийн мод ба бүтэцтэй хайлт}

Энгийн RAG бүх баримтын хэсгийг нэг жагсаалт мэт хайдаг. Сүүлийн үеийн аргууд мэдлэгээ мод эсвэл граф хэлбэрээр зохион байгуулдаг. RAPTOR нь ойролцоо хэсгүүдийг бүлэглэж, LLM-ээр хураангуйлан доороос дээш мод үүсгэдэг \cite{sarthi2024raptor}\unshih. GraphRAG нь текстээс нэр, холбоог LLM-ээр гаргаж граф үүсгэдэг \cite{edge2024graphrag}\unshih. Эдгээр нь маш том текстэд зориулагдсан бөгөөд мод, графаа LLM-ээр бүтээдэг.

Манай мэдлэгийн мод хүний гараар хийсэн JSON файл юм. Баримтын сан 200 орчим хэсэгтэй жижиг. Иймээс модыг хураангуй бичихэд бус, дараах дөрвөн зорилгоор ашиглана (§\ref{sec:knowledge-tree}):
\begin{enumerate}[nosep]
    \item асуултыг тохирох сэдэв рүү чиглүүлэх;
    \item тухайн сэдвийн хэсгүүдийг хайлтад дээш эрэмбэлэх;
    \item багш, өрөө, анги, хуваарийн хүснэгтийг сэдэвтэй холбож шууд хариулах;
    \item мэдээлэл байхгүй сэдвийг таньж, «мэдээлэл алга» гэж хэлэх.
\end{enumerate}
\bichne{W09-д RAPTOR, GraphRAG-ийг уншсаны дараа энэ харьцуулалтыг гүнзгийрүүлнэ.}

\subsection{Хариугүй асуулт}
\bichne{W03 --- \cite{rajpurkar2018squad2}\unshih: сэдэв ойр боловч хариугүй асуулт яагаад хэцүү, систем хэзээ «олдсонгүй» гэх ёстой вэ.}

%======================================================================
\section{Төсөөтэй системүүд}
\bichne{Их сургуулийн элсэлтийн чатботод дүрэм ба RAG-ийг хослуулсан URAG \cite{nguyen2025urag}\unshih зэрэг ажлыг уншиж, манай системтэй хүснэгтээр харьцуулна.}
     % 2. Онол, өнөөгийн түвшин
\chapter{Судалгааны арга зүй}
\label{ch:method}

%======================================================================
% Бүтэц нь яриа таних харьцуулалтын стандарт өгүүлэлтэй ижил:
% өгөгдөл → харьцуулсан систем → туршилтын явц → хэмжүүр → шалгуур.
% Жишээ: Radford нар 2023 (Whisper), Mengke нар 2025, Szymański нар 2020.
% Хайлт (RAG) ба мэдлэгийн модны арга зүйг W07–W08-ын хэмжилтийн дараа
% тусдаа хэсэг болгон нэмнэ. Хуучин текст git түүхэнд (755bf2d) бий.
%======================================================================

Энэ бүлэгт Монгол яриаг текст болгох хоёр системийг хэрхэн харьцуулсныг бичнэ. Судалгааны асуулт нь: \textit{сургуулийн тухай Монгол асуултыг текст болгоход Монгол хэлээр сургасан Whisper загвар ба Gemini 3.5 Transcribe үйлчилгээ ижил бичлэг дээр ямар алдаатай вэ?}

\section{Өгөгдөл}
\label{sec:stt-data}

Нэг хүний 13 бичлэгийг ашигласан. Бичлэг бүр сургуулийн тухай нэг асуулт. Жишээ нь элсэлт, төлбөр, сургуулийн түүхийн тухай асуулт байна. Бичлэгүүдийн нийт урт 152 секунд. Файлууд 44.1 kHz, нэг сувагтай, 16 битийн WAV.

Бичлэг бүрийн зөв бичвэрийг (reference) хүн сонсож баталсан. Зөв бичвэрт нийт 186 үг байна. Хамгийн богино бичлэг 8 үг, хамгийн урт нь 40 үгтэй.

Эдгээр бичлэгээр ямар ч загвар сургаагүй. Зөвхөн шалгахад ашигласан. Гэхдээ бичлэгүүдийг өмнө нь системийг хөгжүүлэхэд харсан. Иймээс энэ нь хөгжүүлэлтийн багц юм, тусдаа шалгалтын багц биш.

\section{Харьцуулсан системүүд}
\label{sec:stt-systems}

Хоёр системд яг ижил аудио файлыг өгсөн. Тохиргоог \tref{tab:stt_setup}-т харуулав.

\begin{table}[H]
\centering
\caption{Харьцуулсан хоёр систем}
\label{tab:stt_setup}
\begin{tabular}{|p{3.0cm}|p{5.4cm}|p{5.2cm}|}
\hline
 & \textbf{Whisper} & \textbf{Gemini 3.5 Transcribe} \\
\hline
Загвар & \texttt{Cafet/\allowbreak whisper-small-\allowbreak mongolian} \cite{cafet_whisper_small_mn} & \texttt{gemini-3.5-transcribe} \cite{google_transcribe_model} \\
\hline
Хувилбар & Hugging Face revision \texttt{d54c63a} & 2026.10.30-ны API \\
\hline
Ажилласан газар & Компьютерын CPU, faster-whisper 1.2.1, int8 & Google-ийн сервер, үнэгүй квот \\
\hline
Тохиргоо & Хэл \texttt{mn}, beam 5, temperature 0, VAD унтраасан & Хэл \texttt{mn-MN}, үгчилсэн (verbatim) горим \\
\hline
\end{tabular}
\end{table}

Whisper загвар нь OpenAI-ийн Whisper small загварыг Монгол хэлээр нэмэлт сургасан хувилбар. Загварын хуудсанд сургалтын өгөгдөл Common Voice 11 гэж бичсэн. Гэхдээ ямар хэсэг дээр, ямар дүрмээр алдааг тооцсон нь тодорхойгүй. Иймээс тэнд бичсэн алдааны хувийг энд ашиглаагүй.

Gemini-ийн загварыг бид сургах, өөрчлөх боломжгүй. Google загвараа шинэчилж болох тул бид туршилтын огноог бичсэн.

\section{Туршилтын явц}
\label{sec:stt-procedure}

Туршилтыг 2026 оны 10-р сарын 30-нд хийсэн. Бичлэг бүрийг систем бүрд нэг удаа өгсөн. Алдаа гарвал дахин илгээгээгүй. Gemini-д илгээсэн файлыг хариу авсны дараа серверээс устгасан.

13 бичлэгийн нэг нь Gemini-ийн хязгаарын алдаа (HTTP 429) авсан. Энэ нь үнэгүй квотын хязгаараас болсон. Хоёр систем хоёулаа хариу өгсөн 12 бичлэгээр харьцуулсан. Энэ 12 бичлэгт 176 үг байна.

Хугацааг бичлэг бүрээр хэмжсэн. Whisper-ийн хугацаа нь компьютер дээр текст болгосон хугацаа. Gemini-ийн хугацаанд файл илгээх, хариу авах, устгах хугацаа орсон. Иймээс хоёр хугацаа яг ижил зүйлийг хэмжээгүй.

\section{Үнэлгээний хэмжүүр}
\label{sec:stt-metrics}

Гол хэмжүүр нь үгийн алдааны хувь (\ac{WER}), томьёо (\ref{eq:wer}). Мөн тэмдэгтийн алдааны хувийг (\ac{CER}) тооцсон. CER-д хоосон зайг тоолоогүй.

Тооцохоос өмнө хоёр системийн текстэд ижил дүрэм хэрэглэсэн:
\begin{enumerate}
    \item Үсгийн хэлбэрийг Unicode NFC болгосон.
    \item Бүх үсгийг жижиг болгосон.
    \item Цэг, таслал зэрэг тэмдгийг зай болгосон.
    \item Тоог өөрчлөөгүй. Жишээ нь «2026» ба «хоёр мянга хорин зургаа» хоёрыг ялгаатай үг гэж үзсэн.
\end{enumerate}
Сүүлийн дүрэм Gemini-д хатуу. Учир нь Gemini тоог цифрээр бичих нь элбэг. Энэ асуудлыг \ref{sec:res-errors}-д дахин ярина.

Сургуулийн мэдээлэлд нэр, тоо, он чухал. Иймээс бичлэг бүрийн чухал үгсийг (нэр, тоо, он) тусад нь шалгасан.

\section{Урьдчилсан шалгуур}
\label{sec:stt-criteria}

Үр дүнг харахаас өмнө шалгуур тогтоосон. Gemini-г Whisper-ээс илүү гэж үзэхийн тулд гурван нөхцөл хангах ёстой:
\begin{enumerate}
    \item Үгийн алдаа Whisper-ээс дор хаяж 10\% бага байх.
    \item Тэмдэгтийн алдаа муудаагүй байх.
    \item Whisper зөв таньсан нэр, тоо, оныг шинээр алдаагүй байх.
\end{enumerate}

\section{Үр дүнгийн тодорхойгүй байдал}
\label{sec:stt-ci}

Бичлэг цөөн тул үр дүн санамсаргүй байдлаас их хамаарна. Үүнийг харуулахын тулд bootstrap аргаар 95\%-ийн интервал тооцсон. 12 бичлэгээс буцаан тавих замаар 12 бичлэг санамсаргүй сонгож, хоёр системийн WER-ийг тооцно. Үүнийг 10000 удаа давтана. Хоёр систем ижил бичлэгээр сонгогдоно. Тооцоог \texttt{tools/\allowbreak bootstrap\_stt\_1030.py} програмаар хийсэн.

Бүх бичлэг нэг хүнийх. Нэг хүний бичлэгүүд хоорондоо төстэй байдаг тул энэ интервал бодит тодорхойгүй байдлаас нарийн гарна. Өөрөөр хэлбэл бодит интервал үүнээс өргөн байж болно.
% ИШ ТАТАХ: Blockwise bootstrap өгүүллийг (STT15, arXiv 1912.09508) W20-д уншсаны дараа
% энд иш татна. Уншаагүй өгүүллийг иш татахгүй.
    % 3. Арга зүй
\chapter{Судалгааны үр дүн}
\label{ch:results}

\section{Дуу танилтын харьцуулалт}
\label{sec:res-stt}

13 бичлэгийн нэг нь Gemini-д хязгаарын алдаа (HTTP 429) өгсөн. Иймээс хоёр систем хоёулаа хариу өгсөн 12 бичлэгээр харьцуулав. Эдгээрт нийт 176 үг байна. Үр дүнг \tref{tab:stt_results}-т харуулав.

\begin{table}[H]
\centering
\caption{Ижил 12 бичлэгийн үр дүн}
\label{tab:stt_results}
\begin{tabular}{|l|r|r|}
\hline
\textbf{Хэмжүүр} & \textbf{Whisper} & \textbf{Gemini 3.5 Transcribe} \\
\hline
Үгийн алдаа & 67 / 176 & 71 / 176 \\
WER & 38.07\% & 40.34\% \\
\hline
CER & 17.93\% & 24.85\% \\
\hline
Нэг бичлэгийн хугацаа (медиан) & 3.98 сек & 6.96 сек \\
\hline
\end{tabular}
\end{table}

Whisper 67, Gemini 71 үг алдсан. Ялгаа 4 үг тул их биш. Тэмдэгтийн алдаа Gemini-д илүү их гарсан. Мөн нэг бичлэгт хэлсэн «2026» оныг Gemini «2020» гэж таньсан, Whisper зөв таньсан. Сургуулийн мэдээлэл оноос их хамаардаг тул энэ алдаа чухал. Gemini урьдчилсан шалгуурт тэнцээгүй. Gemini-ийн хугацаанд файл илгээх, хариу авах хугацаа орсон.

Бүх бичлэг нэг хүнийх, тоо нь цөөн. Тиймээс аль арга нь ерөнхийдөө илүү гэж хэлэх боломжгүй.

\section{Сонголт}
\label{sec:res-choice}

Одоогоор чатад Gemini 3.5 Transcribe-ийг ашиглаж байна. Учир нь яриаг текст болгох тусгай үйлчилгээ, компьютерт загвар ажиллуулах шаардлагагүй, үнэгүй квоттой. Гэхдээ интернэт, API түлхүүр хэрэгтэй, аудио гадагш явдаг. Whisper нь компьютер дээрээ ажилладаг ч Монгол хэлэнд сайжруулахын тулд сургах хэрэгтэй \cite{mengke2025khalkha}.

Эцсийн сонголтыг шинэ бичлэгээр дахин харьцуулж, удирдагч багштайгаа ярилцаж хийнэ.

% Дараа нэмэх: \section{Хайлт ба хариултын үр дүн} — W07–W08-ын хэмжилтийн дараа
        % 4. Үр дүн
\chapter{Дүгнэлт}
\label{ch:conclusion}

%======================================================================
% 2026.11.09: сургуулийн 5 бүлгийн бүтцээр (өмнөх ШМТК-ийн дипломын хэв
% маяг): гол үр дүн → практик ач холбогдол → хязгаарлалт → цаашдын ажил.
% 11/07-ны «Хэлэлцүүлэг» бүлгийн хязгаарлалтыг энд, тайлбарыг 4.3-т шилжүүлэв
% (Chapter5_Discussion.tex файл хадгалагдсан, main.tex-д орохгүй).
%======================================================================

\section{Судалгааны гол үр дүн}
\label{sec:concl-main}

Энэ ажилд сургуулийн ресепшнд зориулсан, Монгол хэлээр дуугаар болон бичгээр асуухад сургуулийн албан эхээс хариулдаг туслах системийн прототип хийсэн. Системийн хоёр хэсгийг тусад нь туршсан: дуу таних хэсгийг нэг хүний бичлэгээр, хайлт ба хариултын хэсгийг бичгээр асуусан асуултаар. Дуугаар асуугаад хариулт авах бүтэн урсгал, ресепшний хүснэгтийн хайлтыг энэ ноорогт хэмжээгүй.

\textbf{RQ1.} Монгол хэлээр сургасан Whisper small загвар ба Gemini 3.5 Transcribe үйлчилгээг нэг хүний 12 бичлэг дээр харьцуулсан. Whisper 38.07\%, Gemini 40.34\% үгийн алдаатай гарсан. Бичлэг цөөн тул энэ ялгаа тодорхой биш. Gemini-ийн тэмдэгтийн алдаа их гарсан нь гол төлөв оныг цифрээр бичсэнээс болсон, гэхдээ нэг бичлэгт оныг буруу таньсан. Иймээс энэ хэмжилтээр Gemini-г илүү гэж үзээгүй. Системд Gemini-г зориулалт, зардал, нэр томьёо зааж өгөх боломжоор нь сонгосон.

\textbf{RQ2.} Сургуулийн сайтын мэдээллээс BM25 ба E5-аар хайж, дүрмээр эх сонгодог системийг 40 асуултаар шалгасан. Эхтэй 20 асуултын бүгдийн зөв хэсэг хайлтын эхний 3-т орсон. Гэхдээ хариулт 13-т нь л бүрэн гарсан. Эх нь байхгүй 19 асуултын 7-д систем эхгүйгээр хариулсан. Иймээс сул тал нь хайлт биш, хариулт сонгох алхам ба эхгүй асуултыг танихад байна. Шалгасан гурван өөрчлөлт бүгд урьдчилсан дүрмээр тэнцээгүй.

Ресепшний асуултын ихэнхэд одоогийн сан хариулт агуулаагүй. Энэ нь системийн гол бэрхшээл хайлтын арга биш, өрөө, багш, хуваарийн мэдээлэл дутуу байгааг харуулсан.

\section{Судалгааны практик ач холбогдол}
\label{sec:concl-practical}

\begin{enumerate}
    \item Сургуулийн сайтын мэдээллээр бичгээр асуухад эхийн үгээр хариулдаг прототип ажиллаж байна. Хариулт бүрт эх хавсаргадаг. Гэхдээ эх нь байхгүй 19 асуултын 7-д ч систем эх хавсаргаж хариулсан тул эх харуулах нь зөв хариултын баталгаа биш. Хэрэглэгч эхийг шалгадаг эсэх, ажилтны ачаалал буурах эсэхийг хэмжээгүй.
    \item Мэдлэгийн мод сургуулийн ямар мэдээлэл дутуу байгааг тодорхой харуулсан. Үүгээр сургалтын албанаас мэдээлэл хүсэх жагсаалт гаргасан.
    \item Сургалтын албаны хүснэгтийг шалгаж оруулах програм бэлэн. Мэдээлэл ирэхэд системийг дахин сургахгүй, зөвхөн хүснэгтийг шинэчилнэ.
\end{enumerate}

\section{Хязгаарлалт}
\label{sec:limits}

\textbf{Туршилт 1 (дуу таних):}
\begin{enumerate}
    \item \textbf{Нэг хүн.} Бүх бичлэг нэг хүнийх. Өөр хүний хоолой, аялга, ресепшний чимээтэй орчинд үр дүн өөр байж болно.
    \item \textbf{Цөөн бичлэг.} 12 бичлэг, 176 үг. Интервал өргөн бөгөөд бодит тодорхойгүй байдал үүнээс ч их байж болно (\ref{sec:stt-ci}).
    \item \textbf{Хөгжүүлэлтийн багц.} Бичлэгүүдийг өмнө нь системийг хөгжүүлэхэд харсан. Тусдаа шалгалтын багц биш.
    \item \textbf{Нэг удаагийн хэмжилт.} Gemini-д бичлэг бүрийг нэг удаа илгээсэн. Нэг бичлэг хязгаарын алдаагаар хасагдсан.
    \item \textbf{Үйлчилгээний хувилбар.} Google загвараа шинэчилж болно. Энэ үр дүн 2026.10.30-ны үйлчилгээнд л хамаарна.
    \item \textbf{Тооны дүрэм.} Цифр ба үгээр бичсэн тоог ялгаатай гэж тооцсон нь Gemini-д хатуу байж болно.
\end{enumerate}

\textbf{Туршилт 2 (хайлт ба хариулт):}
\begin{enumerate}
    \item \textbf{Цөөн асуулт.} 40 асуулт, сайтын нэг хуулбар, нэг өдрийн хэмжилт. Статистикийн ялгааг батлаагүй.
    \item \textbf{Асуултыг хүн хянаагүй.} Асуултыг AI туслах бичиж, сайтын тексттэй тулгаж түгжсэн. Хүн тусад нь хянаагүй.
    \item \textbf{«Утга гарсан» нь бүрэн зөв гэсэн үг биш.} Хариултад шаардсан утга байгаа эсэхийг шалгасан. Хариултыг хүн уншиж үнэлээгүй.
    \item \textbf{Бичгийн асуулт.} Дуугаар асуухад яриа танилтын алдаа нэмэгдэнэ \cite{sidiropoulos2022asr}. Үүнийг энэ туршилтад хэмжээгүй.
    \item \textbf{Ресепшний мэдээлэлгүй.} Сургалтын албаны хүснэгт туршилтын үед ороогүй тул ресепшний хүснэгтийн хайлтыг хэмжээгүй.
\end{enumerate}

\section{Цаашдын ажил}
\label{sec:future}

\begin{enumerate}
    \item Сургалтын албаны өрөө, багш, анги, хуваарийн мэдээллийг шалгаж оруулж, мэдлэгийн модны хоосон хэсгийг бөглөх.
    \item Ресепшний асуултаар шинэ, хүн хянасан багц бэлдэж хэмжих.
    \item Өөр 2--3 хүний зөвшөөрөлтэй бичлэгээр, ресепшний бодит орчинд дуу танилтыг дахин хэмжих.
    \item Хариулт сонгох алхмыг сайжруулах. Асуух үгийн хувирлыг засах дүрэм (F1--F3) хоёр дахь багц дээр ахиц өгөөгүй тул ижил дүрмийг дахин санал болгохгүй. Өөр аргыг (жишээ нь эхгүй асуултыг тусад нь таних) хүн хянасан шинэ багцаар л шалгана.
    \item Хариултыг дуугаар унших хэсгийг хүнээр сонсуулж үнэлэх.
\end{enumerate}
     %    Дүгнэлт (4-р бүлгийн төгсгөл)

%---------  Ном зүй, талархал, хавсралт
\renewcommand{\bibname}{Ном зүй}
\addtotoc{Ном зүй}
\bibliographystyle{unsrt}
\bibliography{references}

% \addtotoc{Талархал}

\chapter*{Талархал}
\thispagestyle{empty}

% Ажлын төгсгөлд өөрийн үгээр бичнэ.
   % ажлын төгсгөлд
% \addtotoc{Хавсралт}

\chapter*{Хавсралт}
\captionsetup{list=no}

% Хавсралтад шалгах асуултын багц, системийг ажиллуулах заавар зэргийг оруулна.
   % хавсралт хэрэгтэй үед

%---------  Товчлол (A4, 2–4 хуудас; Монгол, англи) — урьдчилсан хамгаалалтын
%           дараа бичнэ (гарын авлага х.9). Загварын файлуудыг тэр үед нэмнэ.

\end{document}
